% Figure: the construction behind the Clausius inequality. An arbitrary cyclic
% device exchanges heat with many reservoirs; each heat element is routed through a
% reversible Carnot engine working between the reservoir and a single reference
% reservoir at T0. The composite exchanges heat only with T0 and produces net work
% in a cycle, which the Kelvin-Planck statement forbids unless that work is
% non-positive, giving the cyclic integral of dQ/T <= 0.
\begin{figure}[htbp]
\centering
\begin{tikzpicture}[font=\scriptsize, >=stealth]
% reference reservoir at T0 along the bottom
\draw[fill=engblue!10, draw=engblue!60] (-0.3,-2.6) rectangle (8.3,-1.9);
\node[engblue] at (4.0,-2.25) {reference reservoir at $T_0$};
% the cyclic device in the centre
\draw[fill=enggray!12, draw=enggray!70] (3.2,1.0) rectangle (5.6,2.4);
\node[enggray, align=center] at (4.4,1.7) {cyclic\\ device};
% three reservoirs feeding the device
\foreach \x/\T/\lab in {0.4/2.9/$T_1$, 4.4/3.6/$T_2$, 8.0/2.9/$T_3$} {
  \draw[fill=engred!12, draw=engred!60] (\x-0.55,\T) rectangle (\x+0.55,\T+0.55);
  \node[engred] at (\x,\T+0.275) {\lab};
}
% Carnot bridges from each reservoir down to T0
\foreach \x in {0.4, 4.4, 8.0} {
  \draw[fill=englime!18, draw=englime!60!black] (\x-0.42,-0.6) rectangle (\x+0.42,0.1);
  \node[engblack] at (\x,-0.25) {C};
  % arrow reservoir -> Carnot
  \draw[->, engred!70, thick] (\x,2.78) -- (\x,0.12);
  % arrow Carnot -> reference reservoir
  \draw[->, engblue!70, thick] (\x,-0.62) -- (\x,-1.88);
  % work output of each Carnot bridge feeding the device
  \draw[->, enggray, thick] (\x,-0.25) .. controls (\x+1.0,0.4) and (3.0,1.4) .. (3.18,1.5);
}
% net work out of the device
\draw[->, engorange, very thick] (5.62,1.7) -- (7.4,1.7);
\node[engorange, anchor=west] at (7.42,1.7) {$W_{\mathrm{net}}$};
% labels
\node[engred, anchor=south] at (4.4,4.25) {process reservoirs $T_k$};
\node[englime!50!black, anchor=north] at (4.4,-0.65) {reversible Carnot bridges};
\end{tikzpicture}
\caption[The Clausius-inequality construction]{The composite that proves the
Clausius inequality. An arbitrary device runs a cycle while exchanging heat
$\delta Q_k$ with several reservoirs at temperatures $T_k$ (red). Each heat element
is supplied not directly but through a reversible Carnot bridge (green) that draws
heat from a single reference reservoir at $T_0$ (blue) and delivers $\delta Q_k$ to
the device at $T_k$. A reversible bridge delivering $\delta Q_k$ at $T_k$ must draw
$\delta Q_0 = T_0\,\delta Q_k/T_k$ from the reference reservoir. Summed around one
cycle, the composite (device plus all bridges) returns to its start, exchanges heat
only with the single reservoir $T_0$, and produces net work
$W_{\mathrm{net}} = T_0 \oint \delta Q/T$. The Kelvin-Planck statement forbids net
work in a cycle from a single reservoir, so $W_{\mathrm{net}} \le 0$, which forces
$\oint \delta Q/T \le 0$, the Clausius inequality. Equality holds when every step
is reversible; the strict inequality measures the entropy the real cycle generates.}
\label{fig:vol04ch04:clausius-inequality}
\end{figure}
